%==============================================================================
% Chapitre 3 - Dynamique
%==============================================================================

\section{Introduction}

La dynamique est la branche de la mécanique qui étudie les causes du mouvement.

Alors que la cinématique décrit comment un objet se déplace, la dynamique
cherche à expliquer pourquoi ce mouvement évolue.

La dynamique constitue l'une des bases de la physique moderne et joue un rôle
essentiel dans l'étude de la balistique.

\begin{aretenir}

La cinématique décrit les mouvements.

La dynamique explique leurs causes.

\end{aretenir}

\section{La notion de force}

Dans le langage courant, une force est souvent associée à une poussée ou à une
traction.

En physique, une force représente toute interaction capable de modifier le
mouvement d'un objet.

Une force peut :

\begin{itemize}
\item mettre un objet en mouvement ;
\item ralentir un objet ;
\item modifier sa direction ;
\item le déformer.
\end{itemize}

Une force possède :

\begin{itemize}
\item une intensité ;
\item une direction ;
\item un sens ;
\item un point d'application.
\end{itemize}

\section{Exemple intuitif}

Considérons un ballon posé sur le sol.

Tant qu'aucune force significative ne s'exerce sur lui, il reste immobile.

Si une personne lui donne un coup de pied, une force est appliquée.

Le ballon accélère alors dans la direction de cette force.

Cet exemple simple illustre le lien entre force et variation du mouvement.

\section{Première loi de Newton}

La première loi de Newton est parfois appelée principe d'inertie.

Elle indique qu'un objet conserve son état de repos ou de mouvement rectiligne
uniforme tant qu'aucune force résultante ne s'exerce sur lui.

Cette loi peut sembler évidente, mais elle est en réalité contraire à notre
intuition quotidienne.

Dans la vie courante, les frottements masquent souvent ce comportement.

\begin{exemple}

Une balle glissant sur une surface parfaitement sans frottement continuerait
son mouvement indéfiniment.

Dans le monde réel, les frottements finissent toujours par la ralentir.

\end{exemple}

\section{L'inertie}

L'inertie représente la résistance d'un objet à toute modification de son
mouvement.

Plus un objet possède une masse importante, plus son inertie est élevée.

Il sera donc plus difficile :

\begin{itemize}
\item de le mettre en mouvement ;
\item de le ralentir ;
\item de modifier sa direction.
\end{itemize}

\section{Deuxième loi de Newton}

La deuxième loi de Newton établit un lien entre force, masse et accélération.

Lorsqu'une force est appliquée à un objet, celui-ci subit une accélération.

Plus la force est importante, plus l'accélération est élevée.

Plus la masse est importante, plus l'accélération est faible.

\begin{rigueur}

La relation fondamentale de la dynamique est :

%\begin{formule}
    \begin{equation}
        F = m\,a
        \label{eq:newton2}
    \end{equation}
%\end{formule}

où :

\begin{itemize}
    \item $F$ représente la force ;
    \item $m$ représente la masse ;
    \item $a$ représente l'accélération.
\end{itemize}

\end{rigueur}

\section{Troisième loi de Newton}

La troisième loi est souvent résumée par l'expression :

\begin{quote}
À toute action correspond une réaction égale et opposée.
\end{quote}

Lorsqu'un objet exerce une force sur un second objet, celui-ci exerce une force
de même intensité dans la direction opposée.

\begin{exemple}

Lors du tir :

\begin{itemize}
\item les gaz poussent le projectile vers l'avant ;
\item le projectile exerce une réaction vers l'arrière sur l'arme.
\end{itemize}

Cette réaction participe au recul de l'arme.

\end{exemple}

\section{Forces rencontrées en balistique}

Pendant son vol, un projectile est soumis à plusieurs forces.

Les principales sont :

\begin{itemize}
\item la gravité ;
\item la traînée aérodynamique ;
\item les effets liés au vent ;
\item certaines forces gyroscopiques ;
\item les effets liés à la rotation terrestre.
\end{itemize}

L'importance de ces forces dépend fortement :

\begin{itemize}
\item de la distance de tir ;
\item de la vitesse du projectile ;
\item de l'environnement ;
\item de la précision recherchée.
\end{itemize}

\section{Résultante des forces}

Un projectile est généralement soumis à plusieurs forces simultanément.

L'effet réel observé dépend de leur combinaison.

Cette combinaison est appelée résultante des forces.

C'est cette résultante qui détermine l'accélération du projectile et donc
l'évolution de sa trajectoire.

\section{Application à la balistique}

L'étude balistique consiste principalement à déterminer :

\begin{itemize}
\item quelles forces s'exercent sur le projectile ;
\item comment elles évoluent ;
\item comment elles modifient la trajectoire.
\end{itemize}

Tous les modèles balistiques reposent finalement sur cette idée fondamentale.

\begin{simulation}

Dans un moteur balistique, les forces sont généralement calculées à intervalles
réguliers.

À chaque étape du calcul :

\begin{itemize}
\item les forces sont évaluées ;
\item l'accélération est déterminée ;
\item la vitesse est mise à jour ;
\item la nouvelle position est calculée.
\end{itemize}

La précision de la simulation dépend directement de la qualité de ces calculs.

\end{simulation}

\section{À retenir}

\begin{aretenir}

\begin{itemize}
\item Une force modifie le mouvement d'un objet.
\item La masse caractérise son inertie.
\item Les lois de Newton constituent la base de la mécanique classique.
\item La dynamique explique l'évolution des trajectoires.
\item Toute étude balistique repose sur l'analyse des forces appliquées au projectile.
\end{itemize}

\end{aretenir}

